% titel: Blatt 1, A8: Molmasse aus Isotopenverteilung / A9: Grösste Atomzahl
% Teil II: der durchgerechnete Weg.

\subsection*{A8 --- Molmasse von natürlichem Eisen}
Die Verteilung zeigt in etwa: ${}^{54}$Fe $5.8\%$, ${}^{56}$Fe $91.7\%$,
${}^{57}$Fe $2.2\%$, ${}^{58}$Fe $0.3\%$ (Werte aus der Grafik). Mittlere
Molmasse:
\begin{equation}
  \bar M_{\mathrm{Fe}} = 0.058\cdot54 + 0.917\cdot56 + 0.022\cdot57
  + 0.003\cdot58 = 55.91\ \mathrm{g/mol}.
\end{equation}

\begin{ergebnis}
$M_{\mathrm{Fe}} \approx 55.9\ \mathrm{g/mol}$ (Tabellenwert 55.85; kleine
Abweichung je nach abgelesener Häufigkeit).
\end{ergebnis}

\subsection*{A9 --- Welche Probe hat die meisten Atome?}
\begin{align}
  n_{\mathrm{Fe}}  &= \tfrac{10.0\ \mathrm{g}}{55.8\ \mathrm{g/mol}} = 0.179\ \mathrm{mol},\\
  n_{\mathrm{H_2O}} &= \tfrac{1.0\ \mathrm{g}}{18.0\ \mathrm{g/mol}} = 0.0556\ \mathrm{mol},\\
  n_{\mathrm{He}}  &= \tfrac{1.5\ \mathrm{g}}{4.0\ \mathrm{g/mol}} = 0.375\ \mathrm{mol},\\
  n_{\mathrm{C_6H_6}} &= \tfrac{3.0\ \mathrm{g}}{78.0\ \mathrm{g/mol}} = 0.0385\ \mathrm{mol}.
\end{align}
Helium hat die grösste Stoffmenge und damit die meisten Atome.

\begin{ergebnis}
1.5 g Helium (0.375 mol) --- die meisten Atome, obwohl nicht die grösste Masse.
\end{ergebnis}

\begin{kontrolle}
He ist mit $M=4$ g/mol das leichteste, also hat jede gegebene Masse davon die
grösste Stoffmenge. Reihenfolge $n$: He > Fe > H$_2$O > Benzol.
\end{kontrolle}

\begin{merksatz}
Für atomare Vergleiche nie direkt die Masse vergleichen, sondern $n=m/M$.
Kleinere Molmasse bei gleicher Masse $\Rightarrow$ mehr Stoffmenge/Atome.
\end{merksatz}
